\documentclass[11pt,a4paper]{article}
\usepackage[utf8]{inputenc}
\usepackage[T1]{fontenc}
\usepackage[polish]{babel}
\usepackage{lmodern}
\usepackage[margin=2.3cm]{geometry}
\usepackage{booktabs}
\usepackage{array}
\usepackage{enumitem}
\usepackage{titlesec}
\usepackage{xcolor}
\usepackage{fancyhdr}
\usepackage{hyperref}

\definecolor{navy}{RGB}{27,42,86}
\definecolor{lgray}{RGB}{240,240,244}

\titleformat{\section}{\large\bfseries\color{navy}}{\thesection.}{0.5em}{}
\titleformat{\subsection}{\bfseries\color{navy}}{\thesubsection.}{0.5em}{}

\pagestyle{fancy}
\fancyhf{}
\renewcommand{\headrulewidth}{0.6pt}
\fancyhead[L]{\small Neurofizjologia i pomiar EEG}
\fancyhead[R]{\small Instrukcja laboratoryjna nr 5}
\fancyfoot[C]{\small\thepage}

\newcommand{\infobox}[1]{%
  \par\noindent\colorbox{lgray}{\parbox{\dimexpr\linewidth-2\fboxsep}{\vspace{2pt}#1\vspace{2pt}}}\par\vspace{4pt}}

\title{\vspace{-1.5cm}\textbf{\color{navy}Instrukcja laboratoryjna nr 5}\\[4pt]
\Large Metodologia badań psychofizjologicznych III: pełny proces analizy danych EEG -- od sygnału surowego do wyników}
\author{Neurofizjologia i pomiar EEG\\ Psychologia, studia II stopnia --- specjalność neuropsychologia kliniczna}
\date{Czas trwania zajęć: 120 minut}

\begin{document}
\maketitle
\thispagestyle{fancy}

\section{Cele zajęć}
\begin{itemize}[nosep]
  \item Poznanie pełnego procesu (\emph{pipeline}) przetwarzania danych EEG od rejestracji do gotowych wyników.
  \item Nauka identyfikacji i usuwania typowych artefaktów (mrugnięcia, ruchy gałek ocznych, napięcie mięśniowe).
  \item Zapoznanie z filtrowaniem sygnału i metodą analizy składowych niezależnych (ICA).
  \item Praktyczne przeprowadzenie kompletnej ścieżki obróbki danych na przykładowym zbiorze.
\end{itemize}

\section{Wprowadzenie teoretyczne}

\subsection{Etapy przetwarzania danych EEG}
\begin{enumerate}[nosep]
  \item \textbf{Import i inspekcja wzrokowa} surowego zapisu.
  \item \textbf{Filtrowanie} -- filtr górnoprzepustowy (usuwa dryf wolnofalowy, np. $<$0,1--0,5 Hz), dolnoprzepustowy (ogranicza szum wysokoczęstotliwościowy) i filtr notch (usuwa zakłócenia sieciowe 50 Hz).
  \item \textbf{Ponowne odniesienie (re-referencing)}, np. do średniej z elektrod (\emph{average reference}) lub elektrod wyrostka sutkowatego.
  \item \textbf{Detekcja i korekta artefaktów} -- manualna lub automatyczna (progowanie amplitudy, ICA).
  \item \textbf{Segmentacja} danych na epoki związane ze zdarzeniami.
  \item \textbf{Odrzucanie epok} zawierających artefakty rezydualne (\emph{artifact rejection}).
  \item \textbf{Uśrednianie/agregacja} i analiza statystyczna wyników.
\end{enumerate}

\subsection{Typowe artefakty i ich charakterystyka}
\begin{center}
\begin{tabular}{@{}lll@{}}
\toprule
\textbf{Artefakt} & \textbf{Źródło} & \textbf{Charakterystyka w zapisie}\\
\midrule
Mrugnięcia & Ruch powiek & Duża amplituda, kanały czołowe (Fp1/Fp2)\\
Ruchy gałek ocznych & Aktywność mięśni okoruchowych & Wolna fala, kanały czołowe, przeciwna polaryzacja\\
EMG (napięcie mięśni) & Skurcze mięśni czaszki/karku & Wysoka częstotliwość, nieregularne skoki\\
Artefakt sieciowy & Zakłócenia elektryczne 50 Hz & Wąskopasmowy pik na 50 Hz\\
Artefakt ruchowy & Ruch elektrody/kabla & Nagłe, skokowe zmiany amplitudy\\
\bottomrule
\end{tabular}
\end{center}

\subsection{Analiza składowych niezależnych (ICA)}
\textbf{ICA (Independent Component Analysis)} rozkłada wielokanałowy sygnał EEG na statystycznie niezależne składowe, z których część odpowiada aktywności neuronalnej, a część -- artefaktom (np. gałkoocznym, mięśniowym) o charakterystycznym rozkładzie przestrzennym i widmowym. Identyfikacja i usunięcie składowych artefaktowych, a następnie rekonstrukcja sygnału bez nich, pozwala zachować dane z okresów zawierających mrugnięcia zamiast je odrzucać.

\infobox{\textbf{Zasada dobrej praktyki:} Każdy krok pipeline'u powinien być udokumentowany (parametry filtrów, kryteria odrzucania epok, liczba usuniętych składowych ICA) w celu zapewnienia powtarzalności analizy.}

\section{Sprzęt i oprogramowanie}
Środowisko EEGLAB/MATLAB lub MNE-Python, przykładowy surowy zbiór danych EEG (paradygmat oddball z zajęć 3), skrypt startowy pipeline'u udostępniony przez prowadzącego.

\section{Przebieg zajęć (harmonogram, 120 min)}
\begin{center}
\begin{tabular}{@{}p{2.3cm}p{11.7cm}@{}}
\toprule
\textbf{Czas} & \textbf{Aktywność}\\
\midrule
0--15 min & Wykład: etapy pipeline'u, przegląd typowych artefaktów.\\
15--30 min & Wykład: zasada działania ICA i jej zastosowanie do korekcji artefaktów.\\
30--50 min & Demonstracja: filtrowanie, re-referencing i wizualna inspekcja danych.\\
50--60 min & Przerwa.\\
60--85 min & Ćwiczenie praktyczne: identyfikacja artefaktów i uruchomienie ICA (zadanie 1).\\
85--105 min & Ćwiczenie praktyczne: odrzucanie składowych artefaktowych i rekonstrukcja sygnału, segmentacja epok (zadanie 2).\\
105--120 min & Krytyczna ocena jakości uzyskanych danych, dyskusja, podsumowanie.\\
\bottomrule
\end{tabular}
\end{center}

\section{Zadania do wykonania}
\begin{enumerate}[label=\textbf{Zadanie \arabic*.}, leftmargin=2.7cm]
  \item Na udostępnionych surowych danych zastosujcie filtr górno- i dolnoprzepustowy oraz filtr notch, a następnie wykonajcie wizualną inspekcję i oznaczcie widoczne artefakty (typ, kanał, przybliżony czas wystąpienia).
  \item Uruchomcie dekompozycję ICA, zidentyfikujcie co najmniej dwie składowe artefaktowe (na podstawie topografii i widma) i usuńcie je, rekonstruując oczyszczony sygnał.
  \item Wykonajcie segmentację oczyszczonych danych i porównajcie liczbę odrzuconych epok przed i po korekcji ICA; sformułujcie krótki wniosek o jakości danych.
\end{enumerate}

\section{Kryteria zaliczenia ćwiczenia}
Poprawność identyfikacji artefaktów (30\%), poprawność wykonania i interpretacji ICA (40\%), jakość wniosków końcowych (30\%).

\section{Pytania kontrolne}
\begin{enumerate}[nosep]
  \item Jaką funkcję pełni filtr górnoprzepustowy w przetwarzaniu wstępnym EEG?
  \item Jak odróżnić w zapisie EEG artefakt mrugnięcia od ruchu gałek ocznych?
  \item Na czym polega przewaga korekcji ICA nad prostym odrzucaniem epok z artefaktami?
  \item Dlaczego dokumentowanie parametrów pipeline'u jest istotne metodologicznie?
\end{enumerate}

\section{Literatura}
\begin{enumerate}[nosep]
  \item Chaddad, A., Wu, Y., Kateb, R., \& Bouridane, A. (2023). Electroencephalography signal processing. \emph{Sensors, 23}(14), 6434.
  \item Biasiucci, A., Franceschiello, B., \& Murray, M. M. (2019). Electroencephalography. \emph{Current Biology, 29}(3), R80--R85.
  \item Szumska, I. (2019). \emph{Neurofizjologia i neuropsychologia -- skrypt dla studentów AEH}. Wydawnictwo AEH.
\end{enumerate}

\end{document}
